% Paper template for the <AlgorithmName> RL agent.
%
% This document is based on the arxiv LaTeX style template by George Kour:
%   https://github.com/kourgeorge/arxiv-style (MIT License)
% The MIT license notice (License.txt) is kept with the upstream template
% in rl/template/paper/arxiv-style/. Papers built from this skeleton credit
% the upstream project via the title-page footnote below and the
% kour_arxiv_style entry in references.bib; mention/link the upstream
% project if you redistribute this skeleton, and cite it when the resulting
% paper is published.
%
% The implementation this paper accompanies lives in ../code/. State the
% algorithm's origin, authorship and associated publications below (see the
% repository citation policy in AGENTS.md); do not assume a thesis origin
% unless documented.
%
% Build: pdflatex main && bibtex main && pdflatex main && pdflatex main

\documentclass{article}

\usepackage{arxiv}

\usepackage[utf8]{inputenc} % allow utf-8 input
\usepackage[T1]{fontenc}    % use 8-bit T1 fonts
\usepackage{hyperref}       % hyperlinks
\usepackage{url}            % simple URL typesetting
\usepackage{booktabs}       % professional-quality tables
\usepackage{amsfonts}       % blackboard math symbols
\usepackage{nicefrac}       % compact symbols for 1/2, etc.
\usepackage{microtype}      % microtypography
\usepackage{cleveref}       % smart cross-referencing
\usepackage{graphicx}
\usepackage{natbib}
\usepackage{doi}

\title{<AlgorithmName>\thanks{Typeset using the \emph{arxiv} \LaTeX{} style
	by George Kour, \url{https://github.com/kourgeorge/arxiv-style}
	(MIT license), see \citep{kour_arxiv_style}.}}

% Here you can change the date presented in the paper title
%\date{September 9, 1985}
% Or remove it
%\date{}

\newif\ifuniqueAffiliation
% Comment to use multiple affiliations variant of author block
\uniqueAffiliationtrue

\ifuniqueAffiliation % Standard variant of author block
\author{ \href{https://orcid.org/0000-0000-0000-0000}{\includegraphics[scale=0.06]{orcid.pdf}\hspace{1mm}Author Name}\thanks{Use footnote for providing further
		information about author (webpage, alternative
		address)---\emph{not} for acknowledging funding agencies.} \\
	Department of Computer Science\\
	Cranberry-Lemon University\\
	Pittsburgh, PA 15213 \\
	\texttt{author@university.edu} \\
	%% examples of more authors
	%% \And
	%% Coauthor \\
	%% Affiliation \\
	%% Address \\
	%% \texttt{email} \\
}
\else
% Multiple affiliations variant of author block
\usepackage{authblk}
\renewcommand\Authfont{\bfseries}
\setlength{\affilsep}{0em}
% box is needed for correct spacing with authblk
\newbox{\orcid}\sbox{\orcid}{\includegraphics[scale=0.06]{orcid.pdf}}
\author[1]{%
	\href{https://orcid.org/0000-0000-0000-0000}{\usebox{\orcid}\hspace{1mm}Author Name\thanks{\texttt{author@university.edu}}}%
}
\affil[1]{Department of Computer Science, Cranberry-Lemon University, Pittsburgh, PA 15213}
\fi

% Uncomment to override  the `A preprint' in the header
%\renewcommand{\headeright}{Technical Report}
%\renewcommand{\undertitle}{Technical Report}
\renewcommand{\shorttitle}{<AlgorithmName>}

%%% Add PDF metadata to help others organize their library
\hypersetup{
pdftitle={<AlgorithmName>},
pdfsubject={cs.LG, cs.AI},
pdfauthor={Author Name},
pdfkeywords={reinforcement learning},
}

\begin{document}
\maketitle

\begin{abstract}
	Summary of the algorithm and results.
\end{abstract}


% keywords can be removed
\keywords{reinforcement learning \and Second keyword \and More}


\section{Introduction}
\label{sec:introduction}

Citations use \verb+natbib+, e.g.\ an example upstream reference
\citep{kour2014real, kour2014fast}.


\section{Related Work}
\label{sec:related-work}

Record the algorithm's origin, authorship and associated publications here.


\section{Method}
\label{sec:method}

Cite the repository implementation \texttt{../code/} \citep{ai_ml_research}
and, where the algorithm is thesis-derived, the PhD thesis
\citep{ronnie_stafford_thesis_2020}; entries live in ../citation.bib.

\begin{equation}
	\xi _{ij}(t)=P(x_{t}=i,x_{t+1}=j|y,v,w;\theta)= {\frac {\alpha _{i}(t)a^{w_t}_{ij}\beta _{j}(t+1)b^{v_{t+1}}_{j}(y_{t+1})}{\sum _{i=1}^{N} \sum _{j=1}^{N} \alpha _{i}(t)a^{w_t}_{ij}\beta _{j}(t+1)b^{v_{t+1}}_{j}(y_{t+1})}}
\end{equation}


\section{Experiments}
\label{sec:experiments}

Results produced by \texttt{python -m ai\_ml\_research.rl.<name>.code.experiment}.

\begin{figure}
	\centering
	\fbox{\rule[-.5cm]{4cm}{4cm} \rule[-.5cm]{4cm}{0cm}}
	\caption{Sample figure caption.}
	\label{fig:fig1}
\end{figure}

\begin{table}
	\caption{Sample table title}
	\centering
	\begin{tabular}{lll}
		\toprule
		\multicolumn{2}{c}{Part}                   \\
		\cmidrule(r){1-2}
		Name     & Description     & Size ($\mu$m) \\
		\midrule
		Dendrite & Input terminal  & $\sim$100     \\
		Axon     & Output terminal & $\sim$10      \\
		Soma     & Cell body       & up to $10^6$  \\
		\bottomrule
	\end{tabular}
	\label{tab:table}
\end{table}


\section{Conclusion}
\label{sec:conclusion}


\bibliographystyle{unsrtnat}
\bibliography{references,../citation}  %%% For arXiv submissions, run bibtex locally and inline the generated .bbl, then comment this line out.


%%% Uncomment this section and comment out the \bibliography{references,../citation} line above to use inline references.
% \begin{thebibliography}{1}
%
% 	\bibitem{kour2014real}
% 	George Kour and Raid Saabne.
% 	\newblock Real-time segmentation of on-line handwritten arabic script.
% 	\newblock In {\em Frontiers in Handwriting Recognition (ICFHR), 2014 14th
% 			International Conference on}, pages 417--422. IEEE, 2014.
%
% \end{thebibliography}


\end{document}
